\section{Interpretability}
\label{sec:interp}
%==============================================================================
A common defence of physics modules is interpretability; we test it honestly.
The \ofcv{} gate does not collapse: across training its per-image output
variance grows from $0.154$ (epoch~1) to $0.216$ (epoch~4) and remains
image-dependent, while the \brf{} mean stays near $0.51$ (i.e., \brf{} behaves
as a near-static prior). Qualitatively (Fig.~\ref{fig:interp}), \ofcv{} maps
localise to transparent regions and darken on opaque inclusions in the same
scene. The honest conclusion is a \emph{dissociation}: \ofcv{} maps are
interpretable and non-trivial, yet---per Secs.~\ref{sec:ablation}
and~\ref{sec:cross}---this does not convert into higher \iou. Interpretability
and accuracy are separate axes, and conflating them would overstate the
contribution.

\begin{figure}[t]
\centering
\includegraphics[width=\columnwidth]{fig12_ofcv_interpretability.pdf}
\caption{\ofcv{} maps are image-dependent and localise to transparent regions,
yet do not improve segmentation accuracy (Secs.~\ref{sec:ablation},
\ref{sec:cross}). Panels illustrate the qualitative behaviour.}
\label{fig:interp}
\end{figure}

%==============================================================================
